\subsection{Kernels, Ranges, and Algebraic Projections}
\begin{definition} \label{def:kernel-and-range-linear-map} \hlt{Kernel and Range}\\
For $T\in\LL(\mathcal{V},\mathcal{W})$, define
$\ker(T)\defeq\{\mathbf{v}\in\mathcal{V}:T\mathbf{v}=\mathbf{0}_{\mathcal{W}}\}$ and
$\operatorname{range}(T)\defeq\{T\mathbf{v}:\mathbf{v}\in\mathcal{V}\}$.\\
The kernel is also called the \hlt{nullspace} of $T$.
\end{definition}

\begin{definition} \hlt{Injective and Surjective Linear Maps}\\
A function $T:\mathcal{V}\to\mathcal{W}$ is \hlt{injective} if $T\mathbf{u}=T\mathbf{v}$ implies $\mathbf{u}=\mathbf{v}$. It is \hlt{surjective} if $\operatorname{range}(T)=\mathcal{W}$.
\end{definition}

\begin{proposition} \label{prop:kernel-and-range-are-subspaces} \hlt{Kernel and Range Are Subspaces}\\
If $T \in \LL(\mathcal{V}, \mathcal{W})$, then $\ker(T)$ is a subspace of $\mathcal{V}$ and $\operatorname{range}(T)$ is a subspace of $\mathcal{W}$.
\begin{proof}
Since $T\mathbf{0}_{\mathcal{V}}=\mathbf{0}_{\mathcal{W}}$, the kernel is nonempty.\\
If $\mathbf{u},\mathbf{v}\in\ker(T)$ and $a\in\mathbb{F}$, then $T(a\mathbf{u}+\mathbf{v})=aT\mathbf{u}+T\mathbf{v}=\mathbf{0}_{\mathcal{W}}$, so $a\mathbf{u}+\mathbf{v}\in\ker(T)$.\\
For the range, if $T\mathbf{u},T\mathbf{v}\in\operatorname{range}(T)$ and $a\in\mathbb{F}$, then $aT\mathbf{u}+T\mathbf{v}=T(a\mathbf{u}+\mathbf{v})\in\operatorname{range}(T)$.\\
\hlt{Structural Commentary}: The kernel and range are the two canonical subspaces attached to a linear map.\\
\hlt{Key Insight}: Linearity makes both preimages of zero and images of vectors closed under linear combinations.
\end{proof}
\end{proposition}

\begin{figure}[H]
\centering
\scalebox{0.75}{\begin{tikzpicture}[
  >=Stealth,
  every node/.style={font=\small},
  ambient/.style={draw=gray!35, line width=0.45pt},
  subspace/.style={draw=gray!60, fill=gray!6, line width=0.45pt},
  maparrow/.style={->, gray!65, line width=0.45pt},
  point/.style={circle, fill=gray!70, inner sep=1.4pt}
]
  \coordinate (vzero) at (2.0,1.75);
  \coordinate (wzero) at (7.55,1.75);

  \draw[ambient, fill=white] (0,0) rectangle (4.2,3.6);
  \draw[ambient, fill=gray!15] (5.6,0) rectangle (9.8,3.6);

  \node[anchor=north west] at (0.28,3.42) {$\mathcal{V}$};
  \node[anchor=north west] at (5.88,3.42) {$\mathcal{W}$};

  \draw[subspace]
    plot[smooth cycle, tension=0.75] coordinates {
      (0.78,1.20)
      (0.83,1.85)
      (1.24,2.58)
      (1.98,2.70)
      (2.80,2.25)
      (3.38,1.80)
      (3.02,1.10)
      (2.38,0.76)
      (1.43,0.86)
    };
  \node at (1.85,2.32) {$\ker(T)$};

  \draw[subspace, fill=gray!4]
    plot[smooth cycle, tension=0.72] coordinates {
      (6.24,1.05)
      (6.45,1.72)
      (6.95,2.55)
      (7.60,3.10)
      (8.25,3.02)
      (8.78,2.33)
      (8.94,1.54)
      (8.52,0.90)
      (7.62,0.38)
      (6.72,0.55)
    };
  \node at (7.63,2.34) {$\operatorname{range}(T)$};

  \node[point] at (vzero) {};
  \node[below=4pt] at (vzero) {$\mathbf{0}$};
  \node[point] at (wzero) {};
  \node[below=4pt] at (wzero) {$\mathbf{0}$};

  \draw[maparrow] (vzero) -- (wzero);
  \draw[maparrow] (4.2,3.6) -- (6.62,2.58);
  \draw[maparrow] (4.2,0.0) -- (6.55,1.03);
\end{tikzpicture}
}
\caption{For a linear map $T:\mathcal{V}\to\mathcal{W}$, the range of $T$ may be a proper subspace of $\mathcal{W}$, and the kernel may contain nonzero vectors. The conditions $\operatorname{range}(T)=\mathcal{W}$ and $\ker(T)=\{\mathbf{0}\}$ characterize surjectivity and injectivity, respectively.}
\label{fig:kernel-range-linear-map}
\end{figure}

\begin{definition} \label{def:projection-onto-subspace-along-complement} \hlt{Projection Along a Complement}\\
An operator $P\in\LL(\mathcal{V})$ is an \hlt{algebraic projection}, or \hlt{idempotent operator}, if $P^2=P$.\\
If $\mathcal{V}=\mathcal{U}\oplus\mathcal{W}$, the \hlt{projection onto $\mathcal{U}$ along $\mathcal{W}$} is $\rho_{\mathcal{U},\mathcal{W}}\in\LL(\mathcal{V})$ where $\rho_{\mathcal{U},\mathcal{W}}(\mathbf{u}+\mathbf{w})\defeq\mathbf{u}$ for $\mathbf{u}\in\mathcal{U}$, $\mathbf{w}\in\mathcal{W}$.
\end{definition}

\begin{proposition} \label{prop:algebraic-projection-range-kernel-idempotent} \hlt{Algebraic Projections and Direct-Sum Splittings}\\
For $P\in\LL(\mathcal{V})$, $P$ is an algebraic projection if and only if $\mathcal{V}=\operatorname{range}(P)\oplus\ker(P)$ and $P=\rho_{\operatorname{range}(P),\ker(P)}$.\\
In particular, if $\mathcal{V}=\mathcal{U}\oplus\mathcal{W}$ and $\rho=\rho_{\mathcal{U},\mathcal{W}}$, then $\rho$ is an algebraic projection, $\operatorname{range}(\rho)=\mathcal{U}$, and $\ker(\rho)=\mathcal{W}$.
\begin{proof}
\begin{enumerate}[label=\roman*.]
\setlength{\itemsep}{0pt}
\item[$(\Rightarrow)$] Suppose $P$ is an algebraic projection, so $P^2=P$. For every $\mathbf{v}\in\mathcal{V}$, $\mathbf{v}=P\mathbf{v}+(\mathbf{v}-P\mathbf{v})$, with $P\mathbf{v}\in\operatorname{range}(P)$ and $\mathbf{v}-P\mathbf{v}\in\ker(P)$. If $\mathbf{x}\in\operatorname{range}(P)\cap\ker(P)$, then $\mathbf{x}=P\mathbf{y}$ and $\mathbf{0}=P\mathbf{x}=P^2\mathbf{y}=\mathbf{x}$. Thus the sum is direct. Also $P$ fixes $\operatorname{range}(P)$ and kills $\ker(P)$, so $P=\rho_{\operatorname{range}(P),\ker(P)}$.
\item[$(\Leftarrow)$] Assume the right-hand side. For $\mathbf{v}=\mathbf{x}+\mathbf{y}$ with $\mathbf{x}\in\operatorname{range}(P)$ and $\mathbf{y}\in\ker(P)$, $P^2\mathbf{v}=P\mathbf{x}=\mathbf{x}=P\mathbf{v}$.\\
Thus $P$ is an algebraic projection. For $\rho_{\mathcal{U},\mathcal{W}}$, the formula $\rho_{\mathcal{U},\mathcal{W}}(\mathbf{u}+\mathbf{w})=\mathbf{u}$ gives range $\mathcal{U}$ and kernel $\mathcal{W}$, and the equivalence above shows that $\rho$ is an algebraic projection.
\end{enumerate}
\hlt{Structural Commentary}: Idempotents keep one summand and kill its complement.\\
\hlt{Key Insight}: Direct-sum coordinates turn a projection into $I\oplus0$.
\end{proof}
\end{proposition}

\begin{proposition} \label{prop:complementary-algebraic-projections} \hlt{Complementary Algebraic Projections}\\
If $\mathcal{V}=\mathcal{U}\oplus\mathcal{W}$, then $\rho_{\mathcal{U},\mathcal{W}}+\rho_{\mathcal{W},\mathcal{U}}=I_{\mathcal{V}}$, and $\rho_{\mathcal{U},\mathcal{W}}\rho_{\mathcal{W},\mathcal{U}} =\rho_{\mathcal{W},\mathcal{U}}\rho_{\mathcal{U},\mathcal{W}}=0$.
\begin{proof}
For $\mathbf{v}=\mathbf{u}+\mathbf{w}$, $\mathbf{u}\in\mathcal{U}$, $\mathbf{w}\in\mathcal{W}$, $(\rho_{\mathcal{U},\mathcal{W}}+\rho_{\mathcal{W},\mathcal{U}})\mathbf{v} =\mathbf{u}+\mathbf{w} =\mathbf{v}$,\\
and $\rho_{\mathcal{U},\mathcal{W}}\rho_{\mathcal{W},\mathcal{U}}\mathbf{v} =\rho_{\mathcal{U},\mathcal{W}}\mathbf{w} =\mathbf{0}$; $\rho_{\mathcal{W},\mathcal{U}}\rho_{\mathcal{U},\mathcal{W}}\mathbf{v} =\rho_{\mathcal{W},\mathcal{U}}\mathbf{u} =\mathbf{0}$. Thus the stated operator identities hold.\\
\hlt{Structural Commentary}: A direct-sum decomposition gives two complementary idempotents.\\
\hlt{Key Insight}: Projection onto one summand kills the other summand.
\end{proof}
\end{proposition}

\begin{definition} \label{def:algebraically-orthogonal-projections-and-resolutions} \hlt{Orthogonal Projections in the Algebraic Sense}\\
Two algebraic projections $\rho,\sigma\in\LL(\mathcal{V})$ are \hlt{algebraically orthogonal}, written $\rho\perp\sigma$, if $\rho\sigma=\sigma\rho=0$.\\
\hlt{Resolution of identity} on $\mathcal{V}$ is a finite family of algebraic projections $(\rho_i)_{i=1}^{k}$ with $\rho_i\perp\rho_j$ for $i\neq j$ and $\sum\limits_{i=1}^{k}\rho_i=I_{\mathcal{V}}$.
\end{definition}

\begin{theorem} \label{thm:resolution-of-identity-gives-direct-sum} \hlt{Resolution Ranges Form a Direct Sum}\\
If $(\rho_i)_{i=1}^{k}$ is a resolution of $I_{\mathcal{V}}$ and $R_i=\operatorname{range}(\rho_i)$, then $\mathcal{V}=\bigoplus\limits_{i=1}^{k}R_i$, and $\ker(\rho_i)=\bigoplus\limits_{j\neq i}R_j$.
\begin{proof}
Since $\sum\limits_{i=1}^{k}\rho_i=I_{\mathcal{V}}$, $\mathbf{v}=\sum\limits_{i=1}^{k}\rho_i\mathbf{v}\in\sum\limits_{i=1}^{k}R_i$.\\
If $\sum\limits_{i=1}^{k}\mathbf{x}_i=\mathbf{0}$ with $\mathbf{x}_i\in R_i$, then for each $1\leq m\leq k$, $\mathbf{0}=\rho_m\sum\limits_{i=1}^{k}\mathbf{x}_i=\mathbf{x}_m$, as $\rho_m$ fixes $R_m$ and kills $R_i$ for $i\neq m$.\\
Thus $\mathcal{V}=\bigoplus\limits_{i=1}^{k}R_i$. Also, for $\mathbf{v}=\sum\limits_{j=1}^{k}\mathbf{x}_j$ with $\mathbf{x}_j\in R_j$, $\rho_i\mathbf{v}=\mathbf{x}_i$, so $\ker(\rho_i)=\bigoplus\limits_{j\neq i}R_j$.\\
\hlt{Structural Commentary}: A resolution decomposes the identity into independent idempotent pieces.\\
\hlt{Key Insight}: Pairwise zero products are the operator form of directness.
\end{proof}
\end{theorem}

\begin{theorem} \label{thm:direct-sum-gives-resolution-of-identity} \hlt{Direct Sums Give Resolutions}\\
If $\mathcal{V}=\bigoplus_{i=1}^{k}\mathcal{U}_i$ and $\rho_i=\rho_{\mathcal{U}_i,\bigoplus_{j\neq i}\mathcal{U}_j}$, then $(\rho_i)_{i=1}^{k}$ is a resolution of $I_{\mathcal{V}}$.
\begin{proof}
For $\mathbf{v}=\sum_{i=1}^{k}\mathbf{u}_i$, $\mathbf{u}_i\in\mathcal{U}_i$, $\sum\limits_{i=1}^{k}\rho_i\mathbf{v} =\sum\limits_{i=1}^{k}\mathbf{u}_i =\mathbf{v}$. Thus $\sum\limits_{i=1}^{k}\rho_i=I_{\mathcal{V}}$.\\
If $i\neq j$, then $\rho_i\rho_j\mathbf{v}=\rho_i\mathbf{u}_j=\mathbf{0}$ and $\rho_j\rho_i\mathbf{v}=\rho_j\mathbf{u}_i=\mathbf{0}$, so $\rho_i\perp\rho_j$.\\
\hlt{Structural Commentary}: Direct-sum coordinates produce mutually annihilating coordinate projections.\\
\hlt{Key Insight}: Coordinate projections add to the identity and kill all other summands.
\end{proof}
\end{theorem}

\begin{theorem} \label{thm:algebra-of-projections} \hlt{Algebra of Projections}\\
Let $\rho,\sigma\in\LL(\mathcal{V})$ be algebraic projections and $\operatorname{char}(\mathbb{F})\neq2$.
\begin{enumerate}[label=(\roman*)]
\setlength{\itemsep}{0pt}
\item $\rho+\sigma$ is an algebraic projection $\Longleftrightarrow \rho\perp\sigma$.\\
Then $\operatorname{range}(\rho+\sigma)=\operatorname{range}(\rho)\oplus\operatorname{range}(\sigma)$; $\ker(\rho+\sigma)=\ker(\rho)\cap\ker(\sigma)$
\item $\rho-\sigma$ is an algebraic projection $\Longleftrightarrow \rho\sigma=\sigma\rho=\sigma$.\\
Then $\operatorname{range}(\rho-\sigma)=\operatorname{range}(\rho)\cap\ker(\sigma)$; $\ker(\rho-\sigma)=\ker(\rho)\oplus\operatorname{range}(\sigma)$.
\item If $\rho\sigma=\sigma\rho$, then $\rho\sigma$ is algebraic projection and $\operatorname{range}(\rho\sigma)=\operatorname{range}(\rho)\cap\operatorname{range}(\sigma)$; $\ker(\rho\sigma)=\ker(\rho)+\ker(\sigma)$.
\end{enumerate}
\begin{proof}
\begin{enumerate}[label=(\roman*)]
\setlength{\itemsep}{0pt}
\item Put $A=\rho\sigma$, $B=\sigma\rho$, and $C=\rho\sigma\rho$. Since $(\rho+\sigma)^2=\rho+\sigma \Longleftrightarrow A+B=0$,\\
$A+B=0$ gives $A+C=\rho(A+B)=0$ and $C+B=(A+B)\rho=0$, hence $B=A$ and $2A=0$.\\
Thus $A=B=0$. The converse is immediate. Assume $\rho\perp\sigma$. Then $\operatorname{range}(\rho+\sigma)\subseteq\operatorname{range}(\rho)+\operatorname{range}(\sigma)$, while $(\rho+\sigma)(\mathbf{x}+\mathbf{y})=\mathbf{x}+\mathbf{y}$ for $\mathbf{x}\in\operatorname{range}(\rho)$, $\mathbf{y}\in\operatorname{range}(\sigma)$. Also $\operatorname{range}(\rho)\cap\operatorname{range}(\sigma)=\{\mathbf{0}\}$, since $\mathbf{z}\in\operatorname{range}(\rho)\cap\operatorname{range}(\sigma)$ and $\mathbf{z}=\sigma\mathbf{w}$ imply $\mathbf{z}=\rho\mathbf{z}=\rho\sigma\mathbf{w}=\mathbf{0}$. Finally, $(\rho+\sigma)\mathbf{v}=0 \Longleftrightarrow \rho\mathbf{v}=0=\sigma\mathbf{v}$.
\item Since $P^2=P\Longleftrightarrow (I_{\mathcal{V}}-P)^2=I_{\mathcal{V}}-P$, $\rho-\sigma$ is an algebraic projection $\Longleftrightarrow$ $(I_{\mathcal{V}}-\rho)+\sigma$ is an algebraic projection. By (i), this is equivalent to $(I_{\mathcal{V}}-\rho)\sigma=\sigma(I_{\mathcal{V}}-\rho)=0$, i.e. $\rho\sigma=\sigma\rho=\sigma$.\\
Assume $\rho\sigma=\sigma\rho=\sigma$ and set $Q=\rho-\sigma$. Then $Q\mathbf{v}\in\operatorname{range}(\rho)\cap\ker(\sigma)$; $Q\mathbf{x}=\mathbf{x}$, $\mathbf{x}\in\operatorname{range}(\rho)\cap\ker(\sigma)$, so $\operatorname{range}(Q)=\operatorname{range}(\rho)\cap\ker(\sigma)$. Also $Q$ kills both $\ker(\rho)$ and $\operatorname{range}(\sigma)$, and $\ker(\rho)\cap\operatorname{range}(\sigma)=\{\mathbf{0}\}$.\\
If $\mathbf{v}\in\ker(Q)$, then $\rho\mathbf{v}=\sigma\mathbf{v}$ and $\mathbf{v}=(\mathbf{v}-\rho\mathbf{v})+\sigma\mathbf{v} \in\ker(\rho)\oplus\operatorname{range}(\sigma)$. Hence $\ker(Q)=\ker(\rho)\oplus\operatorname{range}(\sigma)$.
\item If $\rho\sigma=\sigma\rho$, then $(\rho\sigma)^2=\rho^2\sigma^2=\rho\sigma$. Moreover $\mathbf{x}\in\operatorname{range}(\rho\sigma)$ $\Longleftrightarrow$ $\rho\mathbf{x}=\mathbf{x}=\sigma\mathbf{x}$,\\
so $\operatorname{range}(\rho\sigma)=\operatorname{range}(\rho)\cap\operatorname{range}(\sigma)$. Finally, $\ker(\rho)+\ker(\sigma)\subseteq\ker(\rho\sigma)$, and if $\rho\sigma\mathbf{v}=0$,\\
then $\mathbf{v}=(\mathbf{v}-\sigma\mathbf{v})+\sigma\mathbf{v}$; $\mathbf{v}-\sigma\mathbf{v}\in\ker(\sigma)$, $\sigma\mathbf{v}\in\ker(\rho)$. Thus $\ker(\rho\sigma)=\ker(\rho)+\ker(\sigma)$.
\end{enumerate}
\hlt{Structural Commentary}: Operator identities encode range-kernel geometry.\\
\hlt{Key Insight}: Idempotence is a rigid equation; expanding it forces the cross terms to encode the geometry.
\end{proof}
\end{theorem}

\begin{proposition} \label{prop:shift-operators-finitely-nonzero} \hlt{Shift Operators on Finitely Nonzero Sequences}\\
Define $R,S:\mathbb{F}^{(\mathbb{N})}\to\mathbb{F}^{(\mathbb{N})}$ by $R(x_1,x_2,\ldots)=(0,x_1,x_2,\ldots)$, $S(x_1,x_2,\ldots)=(x_2,x_3,\ldots)$. Then $R,S\in\LL(\mathbb{F}^{(\mathbb{N})})$.\\
Also, $\ker(R)=\{\mathbf{0}\}$, $\operatorname{range}(R)=\{(x_1,x_2,\ldots)\in\mathbb{F}^{(\mathbb{N})}:x_1=0\}$, $\ker(S)=\operatorname{span}((1,0,0,\ldots))$, $\operatorname{range}(S)=\mathbb{F}^{(\mathbb{N})}$.
\begin{proof}
Both $R$ and $S$ preserve finite support, so they map $\mathbb{F}^{(\mathbb{N})}$ into itself.\\
Since both maps act componentwise by shifting entries, $R(a\mathbf{x}+\mathbf{y})=aR\mathbf{x}+R\mathbf{y}$ and $S(a\mathbf{x}+\mathbf{y})=aS\mathbf{x}+S\mathbf{y}$ for all $\mathbf{x},\mathbf{y}\in\mathbb{F}^{(\mathbb{N})}$ and $a\in\mathbb{F}$. Thus $R$ and $S$ are linear.\\
The map $R$ has zero kernel as $R\mathbf{x}=\mathbf{0}$ forces every entry of $\mathbf{x}$ to be zero, and its range is exactly the set of finitely nonzero sequences with first entry zero. For $S$, $S\mathbf{x}=\mathbf{0}$ forces $x_2=x_3=\cdots=0$, so $\ker(S)=\operatorname{span}((1,0,0,\ldots))$. Finally, for any $(y_1,y_2,\ldots)\in\mathbb{F}^{(\mathbb{N})}$, we have $S(0,y_1,y_2,\ldots)=(y_1,y_2,\ldots)$, so $S$ is surjective.\\
\hlt{Structural Commentary}: Shift operators expose behavior that finite-dimensional maps cannot have.\\
\hlt{Key Insight}: The right shift is injective but not surjective, while the left shift is surjective but not injective.
\end{proof}
\end{proposition}

\begin{example} \label{ex:shift-operators-show-operator-algebra-noncommutative} \hlt{Shift Operators Show Noncommutativity}\\
For the shifts $R,S\in\LL(\mathbb{F}^{(\mathbb{N})})$ from Proposition \ref{prop:shift-operators-finitely-nonzero}, $SR=I_{\mathbb{F}^{(\mathbb{N})}}$ and $RS(x_1,x_2,\ldots)=(0,x_2,x_3,\ldots)$.\\
Hence $RS\neq SR$, for instance on $(1,0,0,\ldots)$. Therefore $\LL(\mathbb{F}^{(\mathbb{N})})$ is not commutative under composition.
\end{example}
